%"The PDF file may contain up to 25 pages of reference material, single-sided, letter or A4 size, with text and illustrations readable by a person with correctable eyesight without magnification from a distance of 1/2 meter."
\documentclass[10pt,landscape,twocolumn,a4paper,notitlepage]{article}
\usepackage{hyperref}
\usepackage[english, activeacute]{babel}
\usepackage[utf8]{inputenc}
\usepackage{fancyhdr}
\usepackage{lastpage}
\usepackage{listings}
\usepackage{amssymb}
\usepackage[usenames,dvipsnames]{color}
\usepackage{graphicx}
\usepackage{wrapfig}
\usepackage{amsmath}
\usepackage{makeidx}

%%% Margenes
\setlength{\columnsep}{0.25in}    % default=10pt
\setlength{\columnseprule}{0.5pt}    % default=0pt (no line)

\addtolength{\textheight}{2.35in}
\addtolength{\topmargin}{-0.9in}     % ~ -0.5 del incremento anterior

\addtolength{\textwidth}{1.1in}
\addtolength{\oddsidemargin}{-0.55in} % -0.5 del incremento anterior

\setlength{\headsep}{0.08in}
\setlength{\parskip}{0in}
\setlength{\headheight}{15pt}
\setlength{\parindent}{0mm}

%%% Encabezado y pie de pagina
\pagestyle{fancy}
\fancyhead[LO]{\textbf{\title}}
\fancyhead[C]{\leftmark\ -\ \rightmark}
\fancyhead[RO]{Page \thepage\ of \pageref{LastPage}}
\renewcommand{\headrulewidth}{0.4pt}
\fancyfoot{}
\definecolor{darkblue}{rgb}{0,0,0.4}
%%% Configuracion de Listings
\lstloadlanguages{C++}
\lstnewenvironment{code}
	{%\lstset{	numbers=none, frame=lines, basicstyle=\small\ttfamily, }%
	 \csname lst@SetFirstLabel\endcsname}
	{\csname lst@SaveFirstLabel\endcsname}
\lstset{% general command to set parameter(s)
	language=C++, basicstyle=\small\ttfamily, keywordstyle=\slshape,
	emph=[1]{tipo,usa}, emphstyle={[1]\sffamily\bfseries},
	morekeywords={tint,forn,forsn,fore},
	basewidth={0.47em,0.40em},
	columns=fixed, fontadjust, resetmargins, xrightmargin=5pt, xleftmargin=15pt,
	flexiblecolumns=false, tabsize=2, breaklines,	breakatwhitespace=false, extendedchars=true,
	numbers=left, numberstyle=\tiny, stepnumber=1, numbersep=9pt,
	frame=l, framesep=3pt,
    basicstyle=\ttfamily,
    keywordstyle=\color{darkblue}\ttfamily,
    stringstyle=\color{magenta}\ttfamily,
    commentstyle=\color{RedOrange}\ttfamily,
    morecomment=[l][\color{OliveGreen}]{\#}
}

\lstdefinestyle{C++}{
	language=C++, basicstyle=\small\ttfamily, keywordstyle=\slshape,
	emph=[1]{tipo,usa,tipo2}, emphstyle={[1]\sffamily\bfseries},
	morekeywords={tint,forn,forsn,fore},
	basewidth={0.47em,0.40em},
	columns=fixed, fontadjust, resetmargins, xrightmargin=5pt, xleftmargin=15pt,
	flexiblecolumns=false, tabsize=2, breaklines,	breakatwhitespace=false, extendedchars=true,
	numbers=left, numberstyle=\tiny, stepnumber=1, numbersep=9pt,
	frame=l, framesep=3pt,
    basicstyle=\ttfamily,
    keywordstyle=\color{darkblue}\ttfamily,
    stringstyle=\color{magenta}\ttfamily,
    commentstyle=\color{RedOrange}\ttfamily,
    morecomment=[l][\color{OliveGreen}]{\#}
}

%%% Macros
\def\nbtitle#1{\begin{Large}\begin{center}\textbf{#1}\end{center}\end{Large}}
\def\nbsection#1{\section{#1}}
\def\nbsubsection#1{\subsection{#1}}
\def\nbcoment#1{\begin{small}\textbf{#1}\end{small}}
\newcommand{\comb}[2]{\left( \begin{array}{c} #1 \\ #2 \end{array}\right)}
\def\complexity#1{\texorpdfstring{$\mathcal{O}(#1)$}{O(#1)}}
 \newcommand\cppfile[2][]{
\lstinputlisting[style=C++,linerange={#1}]{#2}
}

\begin{document}
\def\title{UACJ - Everyone's Connected}
.\\[0.2cm]
\centering{\LARGE\textbf{Everyone's Connected}} \\[0.5cm]
\centering{\includegraphics[width=5.5cm]{lain.jpg}}
\tableofcontents\newpage

\section{Template}
\cppfile{template.cpp}
\section{Data structures}
\subsection{STL Algorithms}
\begin{justify}
STL stands for Standard Template Library. It is a library that provides several generic 
classes and functions, allowing programmers to manipulate data structures in an easy and efficient way.
The STL provides a range of algorithms which can be used to manipulate data stored in containers. The following 
list shows some of the algorithms provided by the STL and its functions:
\end{justify}
\textbf{Non-Manipulating Algorithms}
\begin{itemize}
\item \textbf{sort(first\_iterator, last\_iterator)} - Sorts the elements in the range [first, last) in ascending order.
\item \textbf{sort(frst\_iterator, last\_iterator, greater\textless int\textgreater())} - Sorts elements inside the vector, in descending order.
\item \textbf{reverse(first\_iterator, last\_iterator)} - Reverses elements inside a vector.
\item \textbf{*max\_element(first\_iterator, last\_iterator)} - Finds the maximum element of a vector.
\item \textbf{*min\_element(first\_iterator, last\_iterator)} - Finds the minimum element of a vector.
\item \textbf{accumulate(first\_iterator, last\_iterator, initial value of sum)} - Summates all the vector elements.
\item \textbf{count(first\_iterator, last\_iterator, x)} - Counts all occurrences 'x' inside a vector.
\item \textbf{find(first\_iterator, last\_iterator, x)} - Returns an iterator to the first occurrence of 'x' in vector and points to last address if the element is not present.
\item \textbf{binary\_search(first\_iterator, last\_iterator, x)} - Tests if 'x' exists in sorted vector or not.
\item \textbf{lower\_bound(first\_iterator, last\_iterator, x)} - Returns an element pointing to the first element in range [first, last), which has a value less than 'x'.
\item \textbf{upper\_bound(first\_iterator, last\_iterator, x)} - Returns an element pointing to the first element in range [first, last), which has a value greater than 'x'.
\end{itemize}
\textbf{Manipulating Algorithms}
\begin{itemize}
\item \textbf{arr.erase(position to delete)} - Erases selected element in vector and shifts and resizes it accordingly.
\item \textbf{arr.erase(unique(arr.begin(), arr.end()), arr.end())} - Erases the duplicate occurrences in sorted vector in a single line.
\item \textbf{next\_permutation(first\_iterator, last\_iterator)} - Modifies the vector to its next permutation.
\item \textbf{prev\_permutation(first\_iterator, last\_iterator)} - Modifies the vector to its previous permutation.
\item \textbf{distance(first\_iterator, desired\_iterator)} - Returns the distance of the desired position from the first iterator to a desired one.
\end{itemize}
\subsection{Binary Search}
\cppfile{Data-Structures/binary_search.cpp}
\subsection{Simplified DSU (Stolen from GGDem)}
\subsection{Disjoint Set Union}
\subsection{Segment Tree}
\subsection{Segment Tree Lazy}
\subsection{Trie}


\section{Graphs}
\subsection{Graph Transversal}
\subsubsection{BFS}
\cppfile{Graphs/breadth_first_search.cpp}
\subsubsection{DFS}
\cppfile{Graphs/depth_first_search.cpp}
\subsection{Topological Sort}
\subsection{APSP: Floyd Warshall}
\subsection{SSSP}
\subsubsection{Lazy Dijkstra}
\subsubsection{Bellman-Ford}
\subsection{Strongly Connected Components: Kosaraju}
\subsection{Articulation Points and Bridges: ModTarjan}

\section{Math}
\subsection{Identities}
{

\textbf{Coeficientes binomiales.}

$(a+b)^{n} = \Sigma_{k = 0}^{n} \binom{n}{k}a^{n-k}b^{k}$

$\binom{n}{k} = \binom{n}{n-k}$

$\binom{n}{k} = \binom{n-1}{k} + \binom{n-1}{k-1}$

$k\binom{n}{k} = n\binom{n-1}{k-1}$

$\Sigma_{k = 0}{n}\binom{n}{k} = 2^{n}$

$\Sigma_{k = 0}^{n} (-1)^{k}\binom{n}{k} = 0$

$\binom{n+m}{t} = \Sigma_{k = 0}^{t}\binom{n}{k}\binom{m}{t-k}$

$\Sigma_{j = k}^{n} \binom{j}{k} = \binom{n+1}{k+1}$

\textbf{Numeros Catalanes.}

$C_n = \frac{2(2n-1)}{n+1} C_{n-1}$

$C_n = \frac{1}{n+1} \binom{2n}{n}$

$C_n \sim \frac{4^n}{n^{3/2}\sqrt{\pi}}$

$\Sigma(n) = O(\log(\log(n)))$ (number of divisors of $n$)

$F_{2n+1} = F_{n}^2 + F_{n+1}^2$

$F_{2n} = F_{n+1}^2 - F_{n-1}^2$

$\sum_{i=1}^n F_i = F_{n+2}-1$

$F_{n+i}F_{n+j} - F_nF_{n+i+j} = (-1)^n F_iF_j$

\textbf{(Möbius Function)}

0 if n is square-free

1 if n got even amount of distinct prime factors

0 if n got odd amount of distinct prime factors

\textbf{(Möbius Inv. Formula)}

Let $g(n) = \sum_{d\mid n} f(d)$, then $f(n)=\sum{d\mid n} g(d) \mu\left(\frac{n}{d})\right)$.
}
\\
\textbf{Permutaciones objetos repetidos}

$P(n,k) = \frac{P(n,k)}{n_{1}!n_{2}!...}$ 
\\
\textbf{Separadores, Ecuaciones lineares $a$ variables $= b$}
  
$\left(\binom{a}{b}\right) = \binom{a+b-1}{b} = \binom{a+b-1}{a-1}$

\textbf{Teorema chino}

sean $\{n_{1}, n_{2},..., n_{k} \}$ primos relativos\\

$P = n_{1}\cdot n_{2}\cdot ...\cdot n_{k}$

$P_{i} = \frac{P}{n_{i}}$

$x\cong a_{1}(n_{1})$

$x\cong a_{2}(n_{2})$
$...$
$x\cong a_{k}(n_{k})$

$P_{1}S_{1} \cong 1(n_{1})$ Donde $S$ soluciones.

$x = P_{1}S_{1}a_{1} + P_{2}S_{2}a_{2} ... P_{k}S_{k}a_{k}$


\subsection{Binary Exponentiation and modArith}
\subsection{Modular Inverse (dividir mod)}
\subsection{Modular Binomial Coeficient and Permutations}
\subsection{Non-Mod Binomial Coeficient and Permutations}
\subsection{Modular Catalan Numbers}
\subsection{Ceil Fraccionario}
\cppfile{Math/ceil_fraccionario.cpp}
\subsection{Numeros de Fibonacci}
\subsection{Sieve Of Eratosthenes}
\subsection{Sieve-based Factorization}
\subsection{Cycle Finding}
\subsection{Berlekamp Massey}
\subsection{Modular Berlekamp Massey}
\subsection{Matrix exponentiation}
\subsection{Ecuaciones Diofantinas}
\subsection{Pollard-Rho, Stolen from GGDem}
\subsection{FFT, Stolen from GGDem}
\subsection{Euler Totient Function}
\section{Geometry}


\section{Strings}
\subsection{Explode by token}
\cppfile{Strings/explode_by_token.cpp}
\subsection{Multiple Hashings DS}
\subsection{Permute chars of string}
\subsection{Longest common subsequence}
\subsection{KMP}
\subsection{Suffix Array}
\subsection{STL Suffix Array}
\section{Classics}
\subsection{Job scheduling}
\subsubsection{One machine, linear penalty}
\subsubsection{One machine, deadlines}
\subsubsection{One machine, profit}
\subsubsection{Two machines, min time}

\section{Flow}
\subsection{Dinic, thx GGDem}

\section{Miscellaneous}
\subsection{pbds}
\subsection{Bit Manipulation}

\section{Testing}
\subsection{Gen and AutoRun testcases}
\subsubsection{Gen.cpp}
\subsubsection{Stress testing}
%\lstinputlisting[language={bash}]{Testing/test.sh}
\subsubsection{Autorun}
%\lstinputlisting[language={bash}]{Testing/runAllTests.sh}
\subsection{Highly Composite Numbers}
Particularly useful when testing number theoretical solutions.
%\lstinputlisting[language={}]{Testing/highlycomposite.txt}


\end{document}

